% vim: set textwidth=78 autoindent:

\section{Uso dei Plugin esterni QGIS Python}\label{sec:external_plugins}\index{plugin}

% when the revision of a section has been finalized, 
% comment out the following line:
%\updatedisclaimer

I plugin esterni QGIS sono scritti in Python e sono mantenuti sia
negli archivi 'Official' e 'User contributed', sia in vari altri archivi esterni 
moderati dai singoli autori. 
La tabella \ref{tab:external_plugins} mostra una lista di plugin attualmente disponibili nell'archivio 'Official', 
con una breve descrizione.
Documentazione dettagliata sull'uso, la minima versione di QGIS richiesta, l'homepage, gli autori 
e altre importanti informazioni vengono fornite insieme con gli stessi plugin esterni 
e non sono incluse in questo manuale.
\footnote{Gli aggiornamenti dei plugin core possono essere
disponibili in questo archivio come fonte esterna.} 
\footnote{I plugin fTools, Mapserver Export e Installer sono anch'essi plugin Python, 
ma sono parte del sorgente QGIS, e sono caricati automaticamente e selezionabili nel gestore dei 
plugin QGIS Plugin Manager (vedere Sezione~\ref{sec:load_external_plugin}).}

Si può trovare una lista aggiornata dei plugin 'Official' nell'archivio ufficiale di QGIS 
in \url{http://qgis.osgeo.org/download/plugins.html}. Questa lista è inoltre disponibile 
automaticamente dal \filename {Plugin installer} via \dropmenuopttwo{plugin_installer}{Scarica Python
Plugins...}. 

\begin{table}[H]
\centering
\caption{Attuali Plugin QGIS esterni moderati}\label{tab:external_plugins}\medskip
\small
 \begin{tabular}{|l|l|p{4in}|}
\hline \textbf{Icona} & \textbf{Plugin esterno} & \textbf{Descrizione}\\
\hline
\includegraphics[width=0.7cm]{zoom2point_icon}
 & Zoom al punto \index{plugin!zoom al punto} & Fa lo zoom alle coordinate 
specificate nella finestra di dialogo di input. Si può specificare il livello di zoom così come 
il controllo dell'estensione della vista.\\
\hline
\includegraphics[width=0.7cm]{plugin_installer}
 & Installatore di Plugin \index{plugins!installatore plugin Python} & La versione più recente dell'Installatore di plugin Python.\\
\hline

\end{tabular}
\end{table}

Una descrizione dettagliata della procedura di installazione per i plugin python esterni si può trovare nella 
Sezione \ref{sec:load_external_plugin}.

\begin{Tip} \caption{\textsc{Aggiungi altri archivi}}
\qgistip{Per aggiungere l'archivio 'User contributed' e/o altri archivi di autori esterni, aprire 
l'Installatore di plugin (\mainmenuopt{plugins} > \dropmenuopttwo{plugin_installer}{Fetch Python Plugins...}),
andare alla scheda \tab{Repositories} e premere \button{Aggiungi repository di terze parti}. 
Se non si vuole uno o più degli archivi aggiunti, disabilitarli con il tasto
\button{Edit...} o rimuoverli completamente con il tasto \button{Delete}.
}
\end{Tip}

